% Generated draft: requires mathematical and source review.
\documentclass[11pt]{article}
\usepackage[margin=25mm]{geometry}
\usepackage{fontspec}
\setmainfont{Latin Modern Roman}
\usepackage{amsmath,amssymb,mathtools,unicode-math}
\setmathfont{Latin Modern Math}
\usepackage{enumitem,needspace,xurl,hyperref}
\setlength{\parindent}{0pt}
\setlength{\parskip}{5pt}
\setlength{\emergencystretch}{3em}
\newcommand{\N}{\mathbb{N}}
\newcommand{\Z}{\mathbb{Z}}
\newcommand{\Q}{\mathbb{Q}}
\newcommand{\R}{\mathbb{R}}
\newcommand{\C}{\mathbb{C}}
\newcommand{\F}{\mathbb{F}}
\newcommand{\A}{\mathbb{A}}
\newcommand{\PP}{\mathbb{P}}
\newcommand{\E}{\mathbb{E}}
\newcommand{\eps}{\varepsilon}
\newcommand{\dd}{\,\mathrm{d}}
\begin{document}
\section*{UnsolvedMath ID: 30003456}
\textbf{Source problem code:} catalogue entry
\subsection{Definitions and mathematical statement}
A global field is either a number field (a finite extension of \(\Q\)) or a global function field (a finite extension of \(\F_q(t)\), where \(q\) is a prime power). The supplied problem concerns automorphic \(L\)-functions over both kinds of global field.

For an automorphic datum \(\pi\) with an associated \(L\)-function \(L(s,\pi)\), let \(s_c\) denote its central point after a normalization has been specified. If \(L(s,\pi)\) is holomorphic near \(s_c\), define its central derivatives by
\[
L^{(r)}(s_c,\pi)=\left.\frac{\mathrm d^r}{\mathrm ds^r}L(s,\pi)\right|_{s=s_c},\qquad r\in\Z_{\geq 0}.
\]
Holomorphy here is a condition needed for this notation to describe finite derivatives, not a hypothesis supplied for a specified family in the record. The record does not specify whether the intended object is an uncompleted, completed, or otherwise normalized \(L\)-function.

\textbf{Research problem.} Develop identities expressing higher-order central derivatives of automorphic \(L\)-functions over number fields and global function fields, extending the Waldspurger and Gross--Zagier formulas in the relevant settings.

\textbf{Unresolved scope.} This is a research program rather than a fully quantified conjecture. The source record does not specify the automorphic groups, representations, precise \(L\)-functions, central-point conventions, derivative orders, local or global hypotheses, or the geometric or arithmetic quantities and normalization constants that should occur in the identities. Consequently, it does not justify an assertion for every automorphic datum or every derivative order. These choices must be supplied before a precise family of identities can be stated.
\subsection{Short English statement}
Seek formulas for higher central derivatives of automorphic \(L\)-functions over number fields and function fields that extend the Waldspurger and Gross--Zagier formulas. The supplied record leaves the families, hypotheses, normalizations, and right-hand sides unspecified, so this draft preserves a research program requiring further mathematical formulation.
\subsection{Sources}
\begin{itemize}
\item Supplied local UnsolvedMath registry record, permanent ID \(30003456\), title ``Higher Central-Derivative Formulas for L-Functions,'' source problem code \texttt{catalogue entry}, in the Oberwolfach Reports Open Problems collection. No individual proposer is identified; the proposed year is supplied as 2017.
\item Supplied report citation: \emph{Arbeitsgemeinschaft: Higher Gross Zagier Formulas (2018)}, Oberwolfach Reports, pp.~1067--1134, DOI \texttt{10.4171/owr/2017/18}. \url{https://doi.org/10.4171/owr/2017/18}. This is also the supplied literature source labeled ``Original Oberwolfach report.'' The citation's 2018 designation and the supplied 2017 provenance require reconciliation.
\item Supplied literature source: \emph{Shtukas and the Taylor expansion of L-functions II}. \url{https://arxiv.org/abs/1712.08026}.
\item Supplied literature assessment dated exactly \texttt{2026-08-21}: the record labels the problem ``partially solved'' and attributes higher central-derivative formulas over function fields to Yun and Zhang, including arbitrary orders within their stated setting; it describes the analogous general number-field geometric theory as incomplete. This is supplied dated curation derived from the OWR collection, not an independently verified literature assessment. Its status and cited claims require independent verification.
\item The supplied statement status is \texttt{corrected\_verified}; its accompanying explanation concerns repairs to notation, typography, and sentence structure while preserving mathematical content. This registry label does not establish a precise conjecture or constitute mathematical review of this draft.
\end{itemize}
\end{document}
